% --- fontspec для LuaLaTeX (заменяет fontenc + inputenc) ---
\usepackage{fontspec}
\usepackage{subcaption}
\usepackage{makecell}
\usepackage{indentfirst}
\usepackage{float}
% --- поля страницы ---
\usepackage[a4paper, margin=2.5cm]{geometry}
% кавычки
\usepackage{csquotes}
% --- babel с LuaLaTeX работает корректно ---
\usepackage[english, russian]{babel}
% --- математика ---
\usepackage{amsmath}
\usepackage{amssymb}
% --- графика (рисунки) ---
\usepackage{graphicx}
% --- таблицы ---
\usepackage{booktabs}
\usepackage{tabularx}
% --- типографика ---
\usepackage[activate={true,nocompatibility},final]{microtype}
% --- графики ---
\usepackage{pgfplots}
\pgfplotsset{compat=1.18}
% --- шрифты ---
\setmainfont{PT Serif}
\setsansfont{PT Sans}
\setmonofont{PT Mono}
% --- библиография (ГОСТ через biblatex+biber) ---
\usepackage[
  backend=biber,
  style=gost-numeric,
  language=auto,
  sorting=none,
]{biblatex}
\addbibresource{inc/bibliography.bib}
% Небольшой запас растяжения для трудно переносимых строк (длинные
% англоязычные названия в списке источников и т.п.).
\setlength{\emergencystretch}{2em}

% --- Титульный блок: УДК, двуязычные аннотации/ключевые слова ---
% Данные кладём в макросы, печатает их переопределённый \@maketitle.
\makeatletter
\newcommand{\udk}[1]{\gdef\@udk{#1}}
\newcommand{\affiliation}[1]{\gdef\@affiliation{#1}}
\newcommand{\email}[1]{\gdef\@email{#1}}
\newcommand{\abstractru}[1]{\gdef\@abstractru{#1}}
\newcommand{\keywordsru}[1]{\gdef\@keywordsru{#1}}
\newcommand{\titleen}[1]{\gdef\@titleen{#1}}
\newcommand{\authoren}[1]{\gdef\@authoren{#1}}
\newcommand{\affiliationen}[1]{\gdef\@affiliationen{#1}}
\newcommand{\abstracten}[1]{\gdef\@abstracten{#1}}
\newcommand{\keywordsen}[1]{\gdef\@keywordsen{#1}}

\renewcommand{\@maketitle}{%
  \noindent УДК~\@udk\par
  \begin{center}
    {\large\bfseries\@title\par}
    \vspace{1em}
    {\bfseries\@author\par}
    \vspace{0.5em}
    \@affiliation\par
    \texttt{\@email}\par
  \end{center}
  \vspace{0.5em}
  {\sloppy
   \noindent\textbf{Аннотация.} \@abstractru\par
   \vspace{0.5em}
   \noindent\textbf{Ключевые слова:} \@keywordsru\par}
  \vspace{1em}
  \begin{center}
    {\large\bfseries\@titleen\par}
    \vspace{1em}
    {\bfseries\@authoren\par}
    \vspace{0.5em}
    \@affiliationen\par
    \texttt{\@email}\par
  \end{center}
  \vspace{0.5em}
  {\sloppy
   \noindent\textbf{Abstract.} \@abstracten\par
   \vspace{0.5em}
   \noindent\textbf{Keywords:} \@keywordsen\par}
  \vspace{1em}
}
\makeatother
